%%%%%%%%%%%%%%%%%%%%%%%%%%%%%%%%%%%%%%%%%
% Jacobs Landscape Poster
% LaTeX Template
% Version 1.0 (29/03/13)
%
% Created by:
% Computational Physics and Biophysics Group, Jacobs University
% https://teamwork.jacobs-university.de:8443/confluence/display/CoPandBiG/LaTeX+Poster
%
% Further modified by:
% Nathaniel Johnston (nathaniel@njohnston.ca)
%
% This template has been downloaded from:
% http://www.LaTeXTemplates.com
%
% License:
% CC BY-NC-SA 3.0 (http://creativecommons.org/licenses/by-nc-sa/3.0/)
%
%%%%%%%%%%%%%%%%%%%%%%%%%%%%%%%%%%%%%%%%%

%----------------------------------------------------------------------------------------
%	PACKAGES AND OTHER DOCUMENT CONFIGURATIONS
%----------------------------------------------------------------------------------------

\documentclass[final]{beamer}

\usepackage[scale=1.24]{beamerposter} % Use the beamerposter package for laying out the poster

\usepackage{amsmath, amssymb, amsthm, mathtools}
\usepackage{graphicx}
\usepackage{cancel}
\usepackage{color}
\usepackage{caption}
\usepackage{subcaption}
%\usepackage{subfigure}
\usepackage{pgf,tikz}
\usetikzlibrary{arrows}

\captionsetup{compatibility=false}

\newcommand{\bs}[1]{\boldsymbol{#1}}
\newcommand{\norm}[1]{\left\| #1 \right\|}
\newcommand{\snorm}[1]{\left| #1 \right|}
\newcommand{\LRp}[1]{\left( #1 \right)}
\newcommand{\LRs}[1]{\left[ #1 \right]}
\newcommand{\LRc}[1]{\left\{ #1 \right\}}
\newcommand{\LRa}[1]{\left\langle #1 \right\rangle}
\newcommand{\LRb}[1]{\left| #1 \right|}
\newcommand{\Grad}{\ensuremath{\nabla}}
\newcommand{\Gradxt}{\ensuremath{\nabla_{xt}}}
\newcommand{\Div}{\ensuremath{\nabla\cdot}}
\newcommand{\Divxt}{\ensuremath{\nabla_{xt}\cdot}}
\newcommand{\Curl}{\ensuremath{\nabla\times}}
\newcommand{\bfH}{\mbox{\boldmath $H$}}
\newcommand{\bfsigma}{\boldsymbol\sigma}
\newcommand{\bfvarsigma}{\boldsymbol\varsigma}
\newcommand{\bftau}{\boldsymbol\tau}
\newcommand{\bfbeta}{\boldsymbol\beta}
\newcommand{\bflambda}{\boldsymbol\lambda}
\newcommand{\bfpsi}{\boldsymbol\psi}
\newcommand{\bfu}{\boldsymbol u}
\newcommand{\bfv}{\boldsymbol v}
\newcommand{\bfV}{\boldsymbol V}
\newcommand{\bfZ}{\boldsymbol Z}
\newcommand{\bfz}{\boldsymbol z}
\newcommand{\bfW}{\boldsymbol W}
\newcommand{\bfw}{\boldsymbol w}
\newcommand{\bfm}{\boldsymbol m}
\newcommand{\bfM}{\boldsymbol M}
\newcommand{\bbM}{\mathbb{M}}
\newcommand{\bfq}{\boldsymbol q}
\newcommand{\bfU}{\boldsymbol U}
\newcommand{\bfS}{\boldsymbol S}
\newcommand{\bbS}{\mathbb{S}}
\newcommand{\bbD}{\mathbb{D}}
\newcommand{\bfK}{\boldsymbol K}
\newcommand{\bbK}{\mathbb{K}}
\newcommand{\bfn}{\boldsymbol n}
\newcommand{\bff}{\boldsymbol f}
\newcommand{\bfF}{\boldsymbol F}
\newcommand{\bbF}{\mathbb{F}}
\newcommand{\bfg}{\boldsymbol g}
\newcommand{\bfG}{\boldsymbol G}
\newcommand{\bfC}{\boldsymbol C}
\newcommand{\bft}{\boldsymbol t}
\newcommand{\bfT}{\boldsymbol T}
\newcommand{\bfI}{\boldsymbol I}
\newcommand{\bbI}{\mathbb{I}}
\newcommand{\bfx}{\boldsymbol x}
\newcommand{\uh}{\widehat{u}}
\newcommand{\fnh}{\widehat{f}_n}
\newcommand{\LQ}{L^2\LRp{Q}}
\newcommand{\LK}{L^2\LRp{K}}
\newcommand{\LVecK}{\mathbf{L}^2\LRp{K}}
\newcommand{\LVecQ}{\mathbf{L}^2\LRp{Q}}
\newcommand{\HdivK}{\bfH(\text{div},K)}
\newcommand{\HdivOmega}{\bfH(\text{div},\Omega)}
% \newcommand{\HdivOmegaLT}{\bfH(\text{div},\Omega)\times L^2([0,T])}
\newcommand{\HdivQ}{\bfH(\text{div}_{xt},Q)}
\newcommand{\HOneK}{H^{1}(K)}
\newcommand{\HOneVecK}{\bfH^{1}(K)}
\newcommand{\HOneQ}{H^{1}(Q)}
\newcommand{\HOneOmegah}{H^{-1}(\Omega_h)}
\newcommand{\HdivOmegah}{\bfH(\text{div},\Omega_h)}
\newcommand{\vdeltau}{v_{\delta\bs u_h}}
\newcommand{\taudeltau}{\bftau_{\delta\bs u_h}}
\newcommand{\ip}[1]{\left\langle #1 \right\rangle}
\newcommand{\pd}[2]{\frac{\partial#1}{\partial#2}}
\newcommand{\pt}[1]{\frac{\partial#1}{\partial t}}
\newcommand{\ppd}[2]{\frac{\partial^2#1}{\partial#2^2}}
\newcommand{\pdd}[3]{\frac{\partial^2#1}{\partial#2\partial#3}}
\newcommand{\der}[2]{\frac{\mathrm{d}#1}{\mathrm{d}#2}}
\newcommand{\Oh}{\Omega_h}
\newcommand{\jump}[1] {\ensuremath{\LRs{\![#1]\!}}}
\newcommand{\Gh}{\Gamma_h}
\newcommand{\mcU}{\mathcal{U}}
\newcommand{\mcUh}{\hat{\mathcal{U}}}
\newcommand{\LOmega}{L^2\LRp{\Omega_h}}

\newcommand{\eqnref}[1]{\eqref{eq:#1}}

\DeclareMathOperator*{\argmin}{arg\,min}
\DeclareMathOperator*{\trace}{tr}

\def\arrtwo#1#2#3#4{\left[
\begin{array}{cc}
#1\; & #2\\
#3\; & #4\\
\end{array}
\right]}
\def\arrthree#1#2#3#4#5#6#7#8#9{\left[
\begin{array}{ccc}
#1\; & #2\; & #3\\
#4\; & #5\; & #6\\
#7\; & #8\; & #9\\
\end{array}
\right]}
\def\arrthreeone#1#2#3{\left[
\begin{array}{ccc}
#1\; & #2\; & #3\\
\end{array}
\right]}
\def\vecttwo#1#2{\left(
\begin{array}{c}
#1\\
#2\\
\end{array}
\right)}
\def\svecttwo#1#2{\left[
\begin{array}{c}
#1\\
#2\\
\end{array}
\right]}
\def\vectthree#1#2#3{\left(
\begin{array}{c}
#1\\
#2\\
#3\\
\end{array}
\right)}
\def\svectthree#1#2#3{\left[
\begin{array}{c}
#1\\
#2\\
#3\\
\end{array}
\right]}

\renewcommand{\arraystretch}{1.2}

\def\etal{{\it et al.~}}

\usetheme{confposter} % Use the confposter theme supplied with this template
\usepackage{exscale}

\setbeamercolor{block title}{fg=utorange,bg=white} % Colors of the block titles
\setbeamercolor{block body}{fg=black,bg=white} % Colors of the body of blocks
\setbeamercolor{block alerted title}{fg=black,bg=utsecgreen} % Colors of the highlighted block titles
\setbeamercolor{block alerted body}{fg=black,bg=dblue!10} % Colors of the body of highlighted blocks

% \setbeamercolor{block title}{fg=ngreen,bg=white} % Colors of the block titles
% \setbeamercolor{block body}{fg=black,bg=white} % Colors of the body of blocks
% \setbeamercolor{block alerted title}{fg=white,bg=dblue!70} % Colors of the highlighted block titles
% \setbeamercolor{block alerted body}{fg=black,bg=dblue!10} % Colors of the body of highlighted blocks
% Many more colors are available for use in beamerthemeconfposter.sty

%-----------------------------------------------------------
% Define the column widths and overall poster size
% To set effective sepwid, onecolwid and twocolwid values, first choose how many columns you want and how much separation you want between columns
% In this template, the separation width chosen is 0.024 of the paper width and a 4-column layout
% onecolwid should therefore be (1-(# of columns+1)*sepwid)/# of columns e.g. (1-(4+1)*0.024)/4 = 0.22
% Set twocolwid to be (2*onecolwid)+sepwid = 0.464
% Set threecolwid to be (3*onecolwid)+2*sepwid = 0.708

\newlength{\sepwid}
\newlength{\onecolwid}
\newlength{\twocolwid}
\newlength{\threecolwid}
\setlength{\paperwidth}{48in} % A0 width: 46.8in
\setlength{\paperheight}{36in} % A0 height: 33.1in
\setlength{\sepwid}{0.024\paperwidth} % Separation width (white space) between columns
\setlength{\onecolwid}{0.22\paperwidth} % Width of one column
\setlength{\twocolwid}{0.464\paperwidth} % Width of two columns
\setlength{\threecolwid}{0.708\paperwidth} % Width of three columns
\setlength{\topmargin}{-0.5in} % Reduce the top margin size
%-----------------------------------------------------------

\usepackage{graphicx}  % Required for including images

\usepackage{booktabs} % Top and bottom rules for tables

%----------------------------------------------------------------------------------------
%	TITLE SECTION
%----------------------------------------------------------------------------------------
\title{Locally Conservative DPG for Fluid Problems} % Poster title

\author{Truman Ellis, Leszek Demkowicz, and Jesse Chan} % Author(s)

\institute{Institute for Computational Engineering and Sciences, University of
Texas} % Institution(s)

%----------------------------------------------------------------------------------------

\begin{document}

\addtobeamertemplate{block end}{}{\vspace*{2ex}} % White space under blocks
\addtobeamertemplate{block alerted end}{}{\vspace*{2ex}} % White space under highlighted (alert) blocks

\setlength{\belowcaptionskip}{2ex} % White space under figures
\setlength\belowdisplayshortskip{2ex} % White space under equations

\begin{frame}[t] % The whole poster is enclosed in one beamer frame

\begin{columns}[t] % The whole poster consists of three major columns, the second of which is split into two columns twice - the [t] option aligns each column's content to the top

\begin{column}{\sepwid}\end{column} % Empty spacer column

\begin{column}{\onecolwid} % The first column

%----------------------------------------------------------------------------------------
%	OBJECTIVES
%----------------------------------------------------------------------------------------

% \begin{alertblock}{Objectives}
%
% Lorem ipsum dolor sit amet, consectetur, nunc tellus pulvinar tortor, commodo eleifend risus arcu sed odio:
% \begin{itemize}
% \item Mollis dignissim, magna augue tincidunt dolor, interdum vestibulum urna
% \item Sed aliquet luctus lectus, eget aliquet leo ullamcorper consequat. Vivamus eros sem, iaculis ut euismod non, sollicitudin vel orci.
% \item Nascetur ridiculus mus.
% \item Euismod non erat. Nam ultricies pellentesque nunc, ultrices volutpat nisl ultrices a.
% \end{itemize}
%
% \end{alertblock}

%----------------------------------------------------------------------------------------
%	INTRODUCTION
%----------------------------------------------------------------------------------------

\begin{block}{Introduction}

The discontinuous Petrov-Galerkin finite element method has been described as least squares finite ele-
ments with a twist. The key difference is that least squares methods seek to minimize the residual of the
solution in the $L^2$ norm, while DPG seeks the minimization in a dual norm realized through the inverse Riesz
map. Exact mass conservation has been an issue that has plagued least squares finite elements for a long
time. In this work, we augment our DPG system with Lagrange multipliers in
order to exactly enforce local conservation. Effectively, this turns our
minimization problem into a constrained minimization problem. We note that
standard DPG, while not guaranteed to be conservative, appears to be nearly
conservative in practice.

\end{block}

%----------------------------------------------------------------------------------------
%	INTRODUCTION
%----------------------------------------------------------------------------------------

\begin{block}{DPG is a Minimum Residual Method}

Let $U$ and $V$ be trial and test Hilbert spaces for a well-posed variational
problem $b(u,v)=l(v)$. In operator form this is $Bu=l$, where $B:U\rightarrow
V'$. We seek to minimize the residual for the discrete space $U_h\subset U$:
\[
u_{h} = \underset{w_{h} \in U_{h}} \argmin \,\, \frac{1}{2}
\norm{Bw_{h}-l}_{V'}^{2}
\]
Use the Riesz inverse to minimize in the $V$-norm rather than its dual:
\begin{align*}
\frac{1}{2}&\norm{Bu_h-l}^2_{V'}=\frac{1}{2}\norm{R_V^{-1}(Bu_h-l)}^2_V\\
&=\frac{1}{2}\LRp{R_V^{-1}(Bu_h-l),R_V^{-1}(Bu_h-l)}_V\,.
\end{align*}
First order optimality requires
the G\^ateaux derivative to be zero in all directions $\delta u \in
U_h$, i.e.,
\[
\left(R_V^{-1}(Bu_h-l),R_V^{-1}B\delta u\right)_V = 0, \quad \forall \delta u \in U_h.
\]
By definition of the Riesz operator, this is equivalent to
\begin{equation*}
\LRa{Bu_h-l,R_V^{-1}B\delta u_h}=0\quad\forall\delta u_h\in U_h\,.
\end{equation*}
Identify $v_{\delta u_h}\coloneqq R_V^{-1}B\delta u_h$ as the
optimal test function for trial function $\delta u_h$. This gives us
\begin{equation*}
b(u_h,v_{\delta u_h})=l(v_{\delta u_h}).
\end{equation*}

\end{block}


%----------------------------------------------------------------------------------------

\end{column} % End of the first column

\begin{column}{\sepwid}\end{column} % Empty spacer column

\begin{column}{\onecolwid} % The second column

%----------------------------------------------------------------------------------------
%	MATERIALS
%----------------------------------------------------------------------------------------

\begin{block}{DPG for Convection-Diffusion}

Start with the strong-form PDE.
\[
\nabla\cdot(\bfbeta u)-\epsilon\Delta u = g
\]
Rewrite as a system of first-order equations.
\begin{align*}
\nabla\cdot(\bfbeta u-\bfsigma)&=g\\
\frac{1}{\epsilon}\bfsigma-\nabla u&=\boldsymbol0
\end{align*}
Multiply by test functions and integrate by parts over each element, $K$.
\begin{align*}
-(\bfbeta u-\bfsigma,\nabla v)_K+((\bfbeta
u-\bfsigma)\cdot\mathbf{n},v)_{\partial K}&=(g,v)_K\\
\frac{1}{\epsilon}(\bfsigma,\bftau)_K+(u,\nabla\cdot\bftau)_K
-(u,\tau_n)_{\partial K}&=0
\end{align*}
Declare traces and fluxes to be independent unknowns and incorporate boundary
conditions to obtain the final variational formulation.
\begin{align*}
-(\bfbeta u-\bfsigma,\nabla v)_K&+(\hat f,v)_{\partial K}
+ \frac{1}{\epsilon}(\bfsigma,\bftau)_K\\
&+(u,\nabla\cdot\bftau)_K
-(\hat u,\tau_n)_{\partial K}=(g,v)_K
\end{align*}

\end{block}

%----------------------------------------------------------------------------------------
%	MATHEMATICAL SECTION
%----------------------------------------------------------------------------------------

\begin{block}{Local Conservation}

The local conservation law in convection diffusion is
\[
\int_{\partial K}\hat f=\int_K g\,,
\]
which is equivalent to having $\mathbf{v}_K:=\{v,\bftau\}=\{1_K,\boldsymbol0\}$ in the test space.
In general, this is not satisfied by the optimal test functions.
Following Moro et al\cite{MoroNguyenPeraire11}, we
can enforce this condition with Lagrange multipliers:
\begin{align*}
L(u_h,\bflambda) = \frac{1}{2}\norm{R_V^{-1}(Bu_h-l)}_V^2
-\sum_K\lambda_K\underbrace{\langle Bu_h-l,\mathbf{v}_K\rangle}_
{\langle\hat f, 1_K\rangle_{\partial K}-\langle g,1_K\rangle_K}\,,
\end{align*}
where $\bflambda=\{\lambda_1,\cdots,\lambda_N\}$.

Finding the critical points of $L(u,\bflambda)$, we get the following
equations.
\begin{align*}
\frac{\partial L(u_h,\bflambda)}{\partial u_h}&=b(u_h,R_V^{-1}B\delta u_h)
-l(R_V^{-1}B\delta u_h)\\
&{\color{red}-\sum_K\lambda_K b(\delta
u_h,\mathbf{v}_K)}=0\quad\forall\delta u_h\in U_h
\end{align*}
\[
\frac{\partial
L(u_h,\bflambda)}{\partial\lambda_K}=-b(u_h,\mathbf{v}_K)+l(\mathbf{v}_K)=0\quad\forall
K
\]

One of that consequences of enforcing local conservation via Lagrange multipliers
is that we've replaced our symmetric positive-definite system with a
saddlepoint problem.

% The robust test norm developed by Chan \etal\cite{ChanHeuerThanhDemkowicz2012}
% includes a $||v||_{L^2}$ with some assumptions that fall apart when the
% convection vector approaches zero. With local conservation, this term is no
% longer needed, and we can replace it with a zero mean term which also results
% in better condition numbers for the local solve.

\end{block}

%----------------------------------------------------------------------------------------

\end{column} % End second column

\begin{column}{\sepwid}\end{column} % Empty spacer column

\begin{column}{\onecolwid} % The third column

%----------------------------------------------------------------------------------------
%	RESULTS
%----------------------------------------------------------------------------------------

\begin{block}{Convection-Diffusion Results}

The locally conservative DPG formulation maintains nearly identical
error convergence behavior as standard DPG for the Erickson-Johnson problem.
\vspace{-1ex}
\begin{figure}
\includegraphics[width=0.8\linewidth]{figs/Erickson/modifiedError.pdf}
% \includegraphics[width=0.48\linewidth]{figs/Erickson/modifiedFlux.pdf}
% \caption*{The locally conservative DPG formulation maintains nearly identical
% error convergence behavior as standard DPG for the Erickson-Johnson problem.}
\end{figure}

% \vspace{3ex}
% Some of the assumptions behind the robust test norm we use break
% down when the convection vector goes to zero, causing standard DPG to overly
% refine parts of the domain with little error. Local conservation allows us to
% avoid these assumptions.
% \begin{figure}
% \centering
% \begin{subfigure}[t]{0.49\textwidth}
% \centering
% \includegraphics[width=1.0\textwidth]{figs/Vortex/modified6nc.png}
% \caption*{\footnotesize{Standard DPG}}
% \end{subfigure}
% \begin{subfigure}[t]{0.49\textwidth}
% \centering
% \includegraphics[width=1.0\textwidth]{figs/Vortex/modified6c.png}
% \caption*{\footnotesize{Locally Conservative DPG}}
% \end{subfigure}
% % \caption*{Some of the assumptions behind the robust test norm we use break
% % down when the convection vector goes to zero, causing standard DPG to overly
% % refine parts of the domain with little error. Local conservation allows us to
% % avoid these assumptions.}
% \end{figure}

\vspace{3ex}
The locally conservative DPG formulation maintains flux imbalances
close to machine precision, even for a discontinuous source term. Standard DPG
becomes more conservative under refinement.
\begin{figure}
\includegraphics[width=0.8\linewidth]{figs/Discontinuous/modifiedFlux.pdf}
% \caption*{The locally conservative DPG formulation maintains flux imbalances
% close to machine precision, even for a discontinuous source term. Standard DPG
% becomes more conservative under refinement.}
\end{figure}

\end{block}

\begin{block}{Inviscid Burgers' Equation Results}

Extension of the theory to other fluid problems like the inviscid Burgers'
equation is fairly trivial.
\begin{figure}
\includegraphics[width=0.8\linewidth]{figs/Burgers/graphFlux.pdf}
% \caption{Inviscid Burgers' equation}
\end{figure}

\end{block}

%----------------------------------------------------------------------------------------

\end{column} % End third column

\begin{column}{\sepwid}\end{column} % Empty spacer column

\begin{column}{\onecolwid} % The fourth column

\begin{block}{Stokes Flow Results}

Least squares finite elements lose 80\% mass flux in Stokes flow
around a cylinder. Locally conservative DPG maintains round-off order flux
imbalance.
\vspace{1ex}
\begin{figure}
\includegraphics[width=0.95\linewidth]{figs/Stokes/graph6c_umag.png}
\caption*{Velocity magnitude}
\end{figure}

\end{block}

%----------------------------------------------------------------------------------------
%	CONCLUSION
%----------------------------------------------------------------------------------------
\vspace{-4ex}

\begin{block}{Conclusions}

\begin{itemize}
  \item We've turned our minimization problem into a saddlepoint problem.
  \item The computational cost is one extra degree of freedom per element.
  % \item Local conservation allows us to eliminate the $||v||_{L^2}$ term in
  %   the inner product.
  \item Enforcement occasionally changes the refinement strategy.
  \item Standard DPG is nearly conservative in practice.
  \item Nearly identical results with better conservation.
  % \item Saw some improvement on the condition number for local solves.
\end{itemize}

\end{block}
%
% %----------------------------------------------------------------------------------------
% %	ADDITIONAL INFORMATION
% %----------------------------------------------------------------------------------------
%
% \begin{block}{Additional Information}
%
% Maecenas ultricies feugiat velit non mattis. Fusce tempus arcu id ligula varius dictum.
% \begin{itemize}
% \item Curabitur pellentesque dignissim
% \item Eu facilisis est tempus quis
% \item Duis porta consequat lorem
% \end{itemize}
%
% \end{block}
%
%----------------------------------------------------------------------------------------
%	REFERENCES
%----------------------------------------------------------------------------------------

\begin{block}{References}

\nocite{*} % Insert publications even if they are not cited in the poster
\small{\bibliographystyle{unsrt}
\bibliography{refs}\vspace{0.75in}}

\end{block}

%----------------------------------------------------------------------------------------
%	ACKNOWLEDGEMENTS
%----------------------------------------------------------------------------------------

% \setbeamercolor{block title}{fg=red,bg=white} % Change the block title color
%
% \begin{block}{Acknowledgements}
%
% \small{\rmfamily{Nam mollis tristique neque eu luctus. Suspendisse rutrum congue nisi sed convallis. Aenean id neque dolor. Pellentesque habitant morbi tristique senectus et netus et malesuada fames ac turpis egestas.}} \\
%
% \end{block}

%----------------------------------------------------------------------------------------
%	CONTACT INFORMATION
%----------------------------------------------------------------------------------------

% \setbeamercolor{block alerted title}{fg=black,bg=utbrown} % Change the alert block title colors
% \setbeamercolor{block alerted body}{fg=black,bg=white} % Change the alert block body colors

% \begin{alertblock}{Contact Information}
%
% \begin{itemize}
% \item Web: \href{http://www.university.edu/smithlab}{http://www.university.edu/smithlab}
% \item Email: \href{mailto:john@smith.com}{john@smith.com}
% \item Phone: +1 (000) 111 1111
% \end{itemize}
%
% \end{alertblock}

\vspace{-6ex}

\begin{center}
\begin{tabular}{ccc}
% \includegraphics[width=0.4\linewidth]{logos/ICES-secondary-logo.png} & \hfill &
\includegraphics[trim=3.0in 7.2in 1.5in 4.0in,width=0.9\linewidth]{logos/ICES-wordmark-teal.pdf}
\end{tabular}
\end{center}

%----------------------------------------------------------------------------------------

%----------------------------------------------------------------------------------------

\end{column} % End of fourth column

%----------------------------------------------------------------------------------------


\begin{column}{\sepwid}\end{column} % Empty spacer column

\end{columns} % End of all the columns in the poster

\end{frame} % End of the enclosing frame

\end{document}
